\section{Related Work}
\label{sec:related}

\subsection{Programming by request}

The ways to program with foundation models have changed constantly over the last
four years, along with the advance in what the models can do. At first a model was
mostly a programming assistant offering inline code completion. Barke et al.\
watched 20 programmers and found them using it in two ways, finishing a line they
had already planned, or asking for a starting point in code they did not yet know
how to write~\cite{barke2023}. Vaithilingam et al.\ found that 24 participants
using Copilot did not complete tasks faster and still preferred to keep it on,
because it gave them a draft to work from~\cite{vaithilingam2022}, and Sarkar et
al.\ collected what the experience felt like from the developer's
side~\cite{sarkar2022}. In all of these the developer accepted a line or a block
of code. Four years later SWE-bench gives a model a natural language issue report
from a real repository and scores it on whether the repository's own tests pass
afterwards~\cite{jimenez2024}, and the systems built to score well on it edit
several files in a loop~\cite{yang2024}. The unit a developer hands over has gone
from a line to a paragraph of English, and the code that comes back is
correspondingly larger.

All of this work studies the minutes around acceptance, asking whether the code
is correct~\cite{perry2023}, whether the developer can tell that it is
correct~\cite{vasconcelos2025}, whether they trust
it~\cite{liang2024,cheng2024}, and how much faster they
finish~\cite{mozannar2024,cui2025}. But did the programmer come to understand the
codebase, understand it enough to explain why it is constructed this way rather
than only to say what it is, understand it enough to recall what their codebase
was a few weeks ago, what has changed since, and what state it is in now? We find
programmers gradually losing the theory of their own
codebase~\cite{naur1985} over the time they use these tools, losing the ability to
describe the rationale behind a design, the ability to carry out a further
modification, and the ability to recall what the codebase held. We suspect this
becomes more of a problem the further specifying in natural language can carry a
developer without their ever having to examine the code.

\subsection{What a version control system agrees to record}

Git records a version as a snapshot of the whole working tree, and each commit
names its parents plus a tree of file contents addressed by hash~\cite{git2026}.
Verifying that a history has not been altered is then a matter of rehashing it,
and a merge is computed on demand from a common ancestor, so nothing in the
repository has to say what happened when two edits met. We keep the first of those
properties and give up the second on purpose. \sgt{} addresses its records by
their contents, and when two edits to one function descend from the same parent it
records that they did rather than computing what the pair amounts to, for the
reason Section~\ref{sec:fork} gives.

Pérez De Rosso and Jackson analysed git's concepts against the purposes users
bring to them, found several places where a concept serves a purpose users do not
have, and removed the staging area and detached \cmd{HEAD} in their
redesign~\cite{perezderosso2013,perezderosso2016}. The commit survives that
redesign untouched, and it should, because it serves a purpose users do have. It
is also the smallest thing git can address, so any operation over less than a
commit works by having the developer point at lines, and \cmd{git add -p} and
\cmd{git rebase -i} are the tools for pointing. Under git's original constraints
this is the cheapest correct design available, because the person running
\cmd{git add -p} in 2005 had typed the lines minutes earlier, and writing their
decomposition into the repository would have meant maintaining a second structure
that could disagree with the first.

Darcs takes the other option and records a version as a patch that adds and
removes lines in a stated context~\cite{roundy2005}. Patches that do not depend on
each other commute, so a user can pull the third patch in a series without the
first two, and the design has been formalised and reimplemented since on a data
structure that makes the operations fast~\cite{jacobson2009,mimram2013,pijul2024}.
A Darcs user who pulls one
patch onto another branch moves the patch that was recorded, along with the
dependencies it names, where git's cherry-pick reapplies a diff in a context it
was not written against and produces a second commit with a second hash, so the
same work has two identities. But Darcs works out what a patch depends on from the
lines the patch touches, so a dependency that leaves no trace in the text goes
unrecorded. When an agent adds a \sym{Retry} class in a new file and edits
\sym{api.py::handle} to call it, the two patches share no line, and nothing in the
record says that pulling the second without the first yields a file referring to a
name that does not exist. Darcs is explicit about this and offers a way to declare
such a dependency by hand, and we would have made the same call in 2003, because
computing a dependency between two functions requires parsing the language and a
system that parses works for some languages and not for others. Darcs declined
that price and we pay it.

\begin{figure*}[t]
  \centering
  % Figure 2. The same hour of work in four records, and what one operation costs
% in each. Same session as Figure 1, same colours, same glyphs.
\newcommand{\FigTwoCaption}{One hour of work with a coding agent, recorded four
ways, with the same request removed from each. The unit a system records decides
which sets a developer can name. Git and Darcs both bottom out at a line, so
removing the caching means selecting lines. Jujutsu and Mercurial add a durable
identity above the line, and that identity belongs to a commit, so a squashed
session carries one identity for all three requests. \sgt{} records one edit per
function together with the request it was made for. The developer supplied that
request and never wrote it down.}

\newcommand{\FigTwoDescription}{A four row comparison. Each row is a card named
after a version control system, drawing the same one hour agent session in that
system's units, with the unit of record stated on the left and the cost of
removing the caching request on the right. The git row draws two commit boxes
filled with coloured diff stripes of four different colours mixed together. The
Darcs and Pijul row draws nine patch tokens in a line with arcs above joining the
patches that share lines, and one dashed arc below joining two patches that depend
on each other through a function call, marked as not recorded. The Jujutsu and
Mercurial row draws one box of mixed stripes, an arrow, and a second box of the
same mixed stripes, labelled with the same change identifier and two different
commit hashes. The sgt row draws four labelled horizontal lines, one per request,
each carrying a mark for every edit filed under that request.}

\begin{tikzpicture}[x=1mm, y=1mm, line cap=round, font=\sffamily]

% ------------------------------- headers ---------------------------------
\node[hd, anchor=west] at (2,79.5)   {UNIT OF RECORD};
\node[hd, anchor=west] at (30,79.5)  {THE SAME HOUR OF WORK, RECORDED FOUR WAYS};
\node[hd, anchor=west] at (117,79.5) {TAKING THE CACHING BACK OUT};

% Every row is one card. The cost of the operation gets a band of its own on
% the right, so the reader can read that column straight down.
% #1 y, #2 height, #3 tab colour, #4 tab text
\newcommand{\cmprow}[4]{%
  \panel{0}{#1}{176}{#2}{#3}{#4}
  \fill[wGrey, rounded corners=0.8mm] (114.5,{#1+1.2}) rectangle (174.6,{#1+#2-1.2});}
% #1 y, #2 text
\newcommand{\rowunit}[2]{%
  \node[anchor=north west, text width=23mm, align=left, inner sep=0, text=sgtInk!75,
        font=\sffamily\fontsize{5.6}{6.5}\selectfont] at (2.2,#1) {#2};}
\newcommand{\rowcost}[2]{%
  \node[anchor=north west, text width=55.5mm, align=left, inner sep=0, text=sgtInk!85,
        font=\sffamily\fontsize{5.7}{6.7}\selectfont] at (116.5,#1) {#2};}

% =============================== git =====================================
\cmprow{55}{19}{sgtGrey}{GIT}
\rowunit{71.4}{a snapshot of the tree. An operation names a line inside a commit.}
\rowcost{71.4}{Select the 6 caching hunks out of 15, in 3 files, across both
  commits, and reverse them. Then read what is left to confirm that nothing
  still calls \fsym{get\_cached}.}

\draw[card, color=sgtInk!18, fill=white] (30,58.2) rectangle (68,70.9);
\node[anchor=north east, inner sep=0, text=sgtGrey,
      font=\ttfamily\fontsize{5}{5.6}\selectfont] at (67.2,70.5) {c39ab2};
\hunk{32.4}{68.5}{9}{sgtGrey}\hunk{32.4}{67.15}{14}{fCache}
\hunk{32.4}{65.8}{12}{fCache}\hunk{32.4}{64.45}{17}{fRate}
\hunk{32.4}{63.1}{16}{fRate}\hunk{32.4}{61.75}{11}{fRate}
\hunk{32.4}{60.4}{19}{fCache}\hunk{32.4}{59.05}{22}{fCache}

\draw[card, color=sgtInk!18, fill=white] (72,58.2) rectangle (110,70.9);
\node[anchor=north east, inner sep=0, text=sgtGrey,
      font=\ttfamily\fontsize{5}{5.6}\selectfont] at (109.2,70.5) {7f10de};
\hunk{74.4}{68.5}{13}{fRetry}\hunk{74.4}{67.15}{18}{fCache}
\hunk{74.4}{65.8}{21}{fRetry}\hunk{74.4}{64.45}{15}{fRetry}
\hunk{74.4}{63.1}{19}{fRetry}\hunk{74.4}{61.75}{10}{fCache}
\hunk{74.4}{60.4}{16}{sgtGrey}

% ========================== Darcs and Pijul ==============================
\cmprow{37}{15.5}{sgtGrey}{DARCS, PIJUL}
\rowunit{50.1}{a patch of added and removed lines, plus the patches whose lines
  it touched.}
\rowcost{50.1}{Unpull the 6 caching patches. \fsym{db.py::query} loses a
  parameter that a caller in another patch still passes, and no recorded
  dependency says so, because the two patches share no line.}

% #1 x centre, #2 colour
\newcommand{\patchtok}[2]{%
  \draw[card, color=#2!85, fill=#2!30, line width=0.35pt]
    ({#1-3.6},44.4) rectangle ({#1+3.6},47.4);}
\patchtok{34}{fRate}\patchtok{42.5}{fRate}\patchtok{51}{fCache}\patchtok{59.5}{fCache}
\patchtok{68}{fCache}\patchtok{76.5}{sgtGrey}\patchtok{85}{fRetry}\patchtok{93.5}{fRetry}
\patchtok{102}{fCache}

\draw[dep] (59.5,47.7) .. controls (63,50.2) and (98.5,50.2) .. (102,47.7);
\draw[dep] (51,47.7) .. controls (54.5,49.7) and (90,49.7) .. (93.5,47.7);
\draw[dep] (68,47.7) .. controls (70,49.2) and (74.5,49.2) .. (76.5,47.7);
\node[acc=sgtInk!70, anchor=south] at (85,50.2) {patches that share lines};
\draw[color=sgtGrey!85, line width=0.35pt, densely dotted]
  (51,44.1) .. controls (56.0,41.3) and (63.0,41.3) .. (68,44.1);
\node[acc=sgtInk!70, anchor=north] at (59.5,42.1)
  {\fsym{handle} calls \fsym{query}. Not recorded.};

% ======================= Jujutsu and Mercurial ===========================
\cmprow{18}{15.5}{sgtGrey}{JUJUTSU, MERCURIAL}
\rowunit{31.1}{a commit whose identity survives being rewritten.}
\rowcost{31.1}{The sitting is one change with one identifier, so the caching has
  no identifier of its own. Split the commit by selecting lines to give the
  caching one.}

\newcommand{\jjbox}[2]{%
  \draw[card, color=sgtInk!18, fill=white] (#1,20.0) rectangle ({#1+30},30.3);
  \hunk{{#1+2.4}}{28.7}{14}{fRate}\hunk{{#1+2.4}}{27.35}{11}{fCache}
  \hunk{{#1+2.4}}{26.0}{19}{fCache}\hunk{{#1+2.4}}{24.65}{16}{fRetry}
  \hunk{{#1+2.4}}{23.3}{22}{fCache}\hunk{{#1+2.4}}{21.95}{12}{sgtGrey}
  \hunk{{#1+2.4}}{20.6}{17}{fRetry}
  \node[anchor=south west, inner sep=0, text=sgtInk!80,
        font=\ttfamily\fontsize{5.2}{6}\selectfont] at (#1,30.7) {#2};}
\jjbox{30}{kxryzmot / 7f10de}
\jjbox{66}{kxryzmot / a19c4b}
\draw[dep] (60.8,25.2) -- (64.6,25.2);
\node[note, anchor=south, text=sgtGrey] at (62.7,25.9) {rewritten};
\annot{98.9}{28.2}{96.6}{30.2}{20}{sgtInk}
\node[anchor=north west, text width=13mm, align=left, inner sep=0, text=sgtInk!70,
      font=\sffamily\fontsize{5.2}{6}\selectfont] at (99.2,29.0)
  {three requests,\\one identifier};

% ================================= sgt ===================================
\cmprow{0}{17}{sgtInk}{\textsf{sgt}}
\rowunit{14.0}{one edit to one function, filed under the request it was made
  for.}
\node[anchor=north west, inner sep=0, text=sgtInk,
      font=\ttfamily\fontsize{5.9}{7}\selectfont] at (116.5,14.0)
  {\$ sgt revert "price cache"};
\node[anchor=north west, text width=55.5mm, align=left, inner sep=0, text=sgtInk!85,
      font=\sffamily\fontsize{5.7}{6.7}\selectfont] at (116.5,10.2)
  {6 edits, 4 functions, 3 files. The 4 edits recorded as built on top of them
   are named before anything is removed.};

\draw[commitcol] (56,2.2) -- (56,14.0);
\draw[commitcol] (86,2.2) -- (86,14.0);
\node[note, text=sgtGrey, anchor=south] at (56,14.1) {c39ab2};
\node[note, text=sgtGrey, anchor=south] at (86,14.1) {7f10de};

\newcommand{\srail}[3]{%
  \draw[rail=#3] (48,#1) -- (112,#1);
  \node[anchor=east, inner sep=0, text=#3,
        font=\sffamily\fontsize{5.8}{6.6}\selectfont] at (46.6,#1) {#2};}
\srail{12.2}{rate limiting}{fRate}
\opadd{56.5}{12.2}{fRate}\opadd{59.9}{12.2}{fRate}\opadd{63.3}{12.2}{fRate}
\opext{66.7}{12.2}{fRate}\ckpt{69.6}{12.2}{fRate}
\srail{9.0}{price cache}{fCache}
\opadd{56.5}{9.0}{fCache}\opadd{59.9}{9.0}{fCache}\opext{63.3}{9.0}{fCache}
\opext{66.7}{9.0}{fCache}\ckpt{69.6}{9.0}{fCache}
\oprew{86.5}{9.0}{fCache}\opext{89.9}{9.0}{fCache}\ckpt{92.8}{9.0}{fCache}
\srail{5.8}{retry policy}{fRetry}
\opadd{86.5}{5.8}{fRetry}\opadd{89.9}{5.8}{fRetry}\opext{93.3}{5.8}{fRetry}
\ckpt{96.2}{5.8}{fRetry}
\srail{2.6}{row accessors}{sgtGrey}
\whoagent{30.0}{2.6}{sgtGrey}
\opmove{56.5}{2.6}{sgtGrey}\opext{59.9}{2.6}{sgtGrey}\opext{63.3}{2.6}{sgtGrey}
\opext{66.7}{2.6}{sgtGrey}\opext{86.5}{2.6}{sgtGrey}

\end{tikzpicture}

  \Description{\FigTwoDescription}
  \caption{\FigTwoCaption}
  \label{fig:granularity}
\end{figure*}

Jujutsu and Mercurial both add a record of history's own history. Jujutsu gives
each change an identifier separate from its commit hash, so rewriting a commit
produces a new commit while the change keeps its name, and a change with more than
one visible commit is called divergent~\cite{jjglossary2026,jujutsu2026}.
Mercurial stores obsolescence markers apart from changeset data, keeping a history
of how the history was edited alongside the history of how the files
were~\cite{mercurial2026}. Both designs let
a developer say ``the change I was working on yesterday'' after five rebases and
be understood, and both let two people rewrite the same region of history and have
the disagreement recorded rather than silently resolved, which is the construction
\sgt{} uses for two competing versions of one function. We part company on what
receives the durable identity. In both systems it is a commit, so the identity is
exactly as coarse as the commit, and after a developer squashes a two-hour agent
session into one commit, one change identifier covers all three requests and
nothing smaller has a name. Figure~\ref{fig:granularity} puts these records beside
each other.

\subsection{Version control that already knew about functions}

Recording a version at the grain of a function has been built before, more than
once, and the reason none of it displaced the commit is the reason we think it can
now. Historage rewrites a Java repository so that every method sits in a file of
its own, at which point git's existing machinery gives each method a history of
its own~\cite{hata2011}, and FinerGit refines that format so that git's rename
detection keeps working when a method moves~\cite{higo2020}. Stellation and Coven
made the fine-grained unit the thing the repository stores rather than a
transformation applied to it, so a developer could check out a single
method~\cite{chucarroll2002,chucarroll2000}. MolhadoRef settled the problem
Section~\ref{sec:identity} reports as our weakest link, by having the development
environment record each refactoring as the developer performed it, so a rename
entered the history as an operation instead of arriving as something a later tool
had to guess at~\cite{dig2007}. Mylyn recorded, for each task a developer worked
on, the set of program elements they touched while working on it, and used that set
to decide what the environment showed them~\cite{kersten2006}. All of this existed
at least fifteen years before the first coding agent and all of it worked.
Historage and FinerGit are in use today as instruments for mining software
histories, which is research reading a repository from the outside, and a working
programmer choosing what to run on Monday morning uses none of it.

We do not read that as a failure of engineering, because a finer grain buys a
developer only what they cannot already supply for themselves, and until recently
they could supply all of it. Whoever wanted the caching out of a 2010 repository
had written the caching, so a system that told them which methods it touched was
performing work they had already done, at the price of a repository their
colleagues could not clone with plain git. What has changed is not the value of
the grain but the identity of the author, and a record that was redundant while
the developer wrote the lines is the only surviving copy of the decomposition once
they have stopped writing them. Mylyn's task context and \sgt{}'s named group of
edits are the same object built for two different situations, and the difference is
that Mylyn watched a developer's attention move through code that developer was
writing, whereas the attention \sgt{} reconstructs belonged to a process that has
already exited.

The merge question in this literature was approached from the end opposite to ours.
Horwitz, Prins and Reps showed that two versions of a program can be integrated
automatically when a semantic analysis proves the changes do not
interfere~\cite{horwitz1989}, and Apel et al.'s semistructured merge parses a file
far enough to see that two edits landed in different methods, resolving a large
share of what a line-based merge reports as a conflict~\cite{apel2011}. The design
here rests directly on two further results from that area, that a merge of the
recorded operations tells a developer more than a merge of the states those
operations produced~\cite{lippe1992}, and that a version is usefully treated as a
set with a logic over it, which is the construction Section~\ref{sec:sets}
uses~\cite{zeller1997}. \sgt{} departs from Horwitz and from Apel deliberately.
Where semistructured merge combines two edits to one
method as soon as it can see they touch different lines, \sgt{} stops and asks a
person, and Section~\ref{sec:forkoff} concedes that a repository with one person in
it is the wrong place to have tested that rule.

\subsection{Recovering a decomposition after the fact}

Commits that address more than one concern distort the tools built from a history.
Herzig and Zeller examined five open source Java projects, found such commits
throughout them, and showed that separating those changes alters which files a
defect prediction model identifies as the most likely to contain
bugs~\cite{herzig2013,herzig2016}. The untangling tools built after that finding
recover the grouping from the diff itself, partitioning it by the definition and
use relationships between the changed statements~\cite{barnett2015} or by a
dependency graph across the versions a commit spans~\cite{partachi2020}, and every
one of them is evaluated on whether it reproduces a grouping a human labelled. Two
anticipate the move we make: Dias et al.\ read the fine-grained edit events an IDE
records while the developer is still typing~\cite{dias2015}, and Shen et al.\ hand
the partition back to the developer to correct, which is the arrangement
Section~\ref{sec:names} adopts~\cite{shen2021,tao2015,wang2019}. What none of them could
assume is a statement of the boundary written down before the code existed, because
in 2015 no such statement existed to be had, and the strength of their inference is
a measure of how much they had to recover without one.

That is the loosened constraint, and it is worth being exact about what it buys,
because \sgt{} performs the same inference they do and performs it no better; we
would expect it to be worse, since it has to finish while a developer waits. What
has changed is what the inference is for. Their partition has to be right, because
it is the answer. Ours has to be close enough that a name attaches to the right
set, and the developer who reads a wrong name fixes it with a rename that moves no
code.

\sgt{} does depend on two results from this literature. Keeping a function's
identity when it is renamed or moved is the problem that refactoring detection and
tree differencing address~\cite{tsantalis2018,falleri2014,fluri2007}, and it is
the weakest link in our implementation, since a rename we fail to detect appears
as one function removed and another added. Grouping edits into features is a
clustering over which functions change together and which call each other, both
halves of which the co-change mining literature
established~\cite{zimmermann2004,ying2004,kim2005,godfrey2005}, resolved with a
community detection algorithm~\cite{traag2019,traag2011}. Both are inferences, and
Section~\ref{sec:design} states what a developer does when they are wrong.

\subsection{Histories the editor keeps}

A second tradition puts the record in the editor and asks nothing of the
developer. Azurite stores a programmer's editing operations at a finer grain than
any commit and lets them undo one past edit while keeping the edits made after
it~\cite{yoon2015,yoon2013,yoon2014}. Variolite lets a data
scientist wrap a region of code as a variant and switch between variants in place,
after Kery et al.\ found that people doing exploratory work avoid version
control for exactly that work~\cite{kery2017}, and the same tradition has held
alternative versions side by side, replayed a document's editing history, versioned
one experiment at a time, and cleaned a notebook up afterwards, in every case
recording what the developer did as a byproduct of their doing it and never asking
them to prepare
it~\cite{hartmann2008,terry2004,grossman2010,mikami2017,kery2018,head2019}. That
commitment \sgt{} keeps. All of them also stop at
the boundary of one machine: the history lives in the editor's own store, so a
colleague, a build, and a code review never see it, and the unit is a single edit
event, so a set of related edits still has no name a person can use. \sgt{} pays
for crossing that boundary by writing its record into the git repository as files
committed with the code, and Section~\ref{sec:colleague} describes what that costs.

\subsection{Layers over git built for the way people work now}

The tools developers reach for today accept git's storage and add a layer above
it. GitButler keeps several branches applied to one working directory at once,
draws each as a vertical lane with its own staging area, and lets the developer
drag a change into the lane it belongs to~\cite{gitbutler2026}, on the correct
reasoning that which branch a piece of work belongs to is easiest to decide once
the work exists. GitButler assigns an individual file change to a lane, so it
cannot express the case at the centre of Figure~\ref{fig:two-records}, where three
requests all edited \sym{api.py::handle} and no assignment of \sym{api.py} to one
lane is right. Graphite and GitHub's stacked pull requests take the ordered series
as their unit, keeping a chain of changes consistent by rebasing and retargeting
the higher layers when a lower one merges and landing every unmerged layer beneath
the topmost ready one in a single
operation~\cite{graphite2026,githubstacked2026}, which is the right answer to a
reviewer who cannot read one large change with the care they can give to each of
five small ones. Graphite's guide states the precondition plainly, that developers
must understand how to break down their work into small, incremental changes, and
every command the tool offers operates on a decomposition the developer has
already performed. In an agent session that decomposition exists in the
developer's request history and nowhere in the code.

The closest system to ours is sem, which parses about 32 languages with
tree-sitter and reports which functions, methods and classes a commit changed,
resolving an entity's identity across versions in three passes, by exact
identifier, by a structural hash that ignores whitespace and comments, and by
token overlap above 80\%~\cite{sem2026}. sem keeps its answers in a derived SQLite
cache outside the
repository, and that placement makes sem safe to adopt and safe to abandon, gives
entity-level answers about a history recorded years before sem existed, and
requires agreement from nobody else on the team. It also bounds what sem can
offer: a record any teammate's ordinary commit can silently invalidate cannot be
the record an operation acts on, so every sem command reports and none of them
changes state, and a developer who learns from \cmd{sem diff} that
\sym{get\_cached} changed removes it with a line-level tool.
\sgt{} commits its record to the repository in
order to act on it, and we pay for that in two places. A record that operations
depend on has to be repaired when it is wrong, which
Section~\ref{sec:unreproducible} reports the current cost of, and everybody who
pulls the repository now carries a record they did not ask for, which
Section~\ref{sec:colleague} works through from the point of view of the three
people who never installed \sgt{}. sem chose the safer placement and we chose the
more useful one, and the honest form of that sentence is that we have taken on an
obligation sem does not have.
